\scanpage{198}
\mspage{II.46}
% archive sheet number: 197
% lines 1--5, with insertion in upper margin
Bien entendu, les exemples 1, 2 \inserted{de rep. de $\Pi^{\unsure{\mathrm{an}}}$ ou $\Pi^{\unsure{\mathrm{alg}}}$} ont un caractère \uline{fonctoriel} pour $X$, resp $A$, variable. En particulier, dans le cas d'une $A$, on trouve un hom d'anneaux (injectif),
\[
\operatorname{End}(A)\otimes_{\mathbb{Z}}\mathbb{Q}_{\unsure{\ell|l}}\hookrightarrow\operatorname{End}_{\Pi}\bigl(\underbrace{T_\ell(A(\unsure{\xi|\eta}))\otimes_{\mathbb{Z}_\ell}\mathbb{Q}_\ell}_{E}\bigr).
\tag*{(\unsure{6.15.1|6.16.1})}
\]

% lines 6--17
\uline{Remarque 6.16.} Une des conséquences des \uline{conjectures de Tate}, (qui montre le rôle essentiel des représentations \unsure{obtenues|linéaires|retenues} des \unsure{gr|gpes|gps} fondamentaux \deleted{\unsure{géom.}} \unsure{arithmétiques|arithmetiques}) est que l'hom précédent est bijectif \deleted{\unsure{Ce}} lorsque $X$ est normal intègre. Pour cette question, on se ramène d'ailleurs aussitôt au cas où $X$ est le spectre d'un corps \uline{algébrique clos}. Cela montre à quel point ($=$ dans un cas soi-disant ``géométrique'') l'opération de $\Pi$ est loin d'être triviale. [C'est prouvé par Tate si $X$ est le spectre d'un corps fini, \deleted{\unsure{et}} ou \deleted{\unsure{partant dans le}} \inserted{de façon équivalente,} si c'est le spectre de la clôture algébrique d'un corps fini .... ].

% lines 18--22
\uline{Remarque} 6.17. On s'est rendu compte que $\Pi$ lui-$=$, ou $G=\Pi^{\mathrm{an}}$, $\unsure{\widetilde G|\overline G}=\Pi^{\mathrm{alg}}$, $\unsure{\mathfrak g^\circ|\mathfrak g^\bullet}$, $\unsure{\widetilde{\mathfrak g}|\overline{\mathfrak g}}$, sont démesurément grands et ils échappent pratiquement à notre contrôle. Leur rôle est surtout théorique, i.e. leurs représentations $\ell$-adiques
